\documentclass[10pt,a4paper]{amsart}
\usepackage[margin=19mm]{geometry}
\usepackage{amsmath,amssymb,mathtools}
\usepackage[scr=rsfs,cal=euler]{mathalfa}
\usepackage{xfrac}
\usepackage[unicode,hidelinks,bookmarks=false]{hyperref}
\newcommand{\ie}{i.e.\ }
\newcommand{\cf}{cf.\ }
\newcommand{\resp}{respectively\ }
\newcommand{\Subsection}{\subsection}
\providecommand{\nobreakdash}{\nobreak\hbox{-}}
\setlength{\emergencystretch}{2em}
\multlinegap0pt
\usepackage{amsthm}
\usepackage{mathtools}
\usepackage{amsmath,amssymb,amsfonts}
\usepackage{graphicx}
\usepackage{bm}
\usepackage{enumitem}
\usepackage{mathrsfs}
\usepackage{xcolor}
\newtheorem{theorem}{Theorem}[section]
\newtheorem{proposition}{Proposition}[section]
\newtheorem{lemma}[theorem]{Lemma}
\newtheorem{remark}[theorem]{Remark}
\newtheorem{example}{Example}[section]
\newtheorem{definition}{Definition}[section]
\newtheorem{corollary}[theorem]{Corollary}
\title[Logvinenko-Sereda for the Dunkl transform]{Observability Inequalities and the Logvinenko-Sereda Theorem for the Dunkl Transform: the uniform cube-measure comparison (Lemma 4.1)}
\author[X. Zhao, L. Wei, Z. Duan]{Xingyu Zhao, Longben Wei, Zhiwen Duan}
\date{Source: arXiv:2609.28993v1 (CC BY 4.0)}
\begin{document}
\maketitle
\vspace{1em}
\noindent\textit{Source and licence.} Observability Inequalities and the Logvinenko-Sereda Theorem for the Dunkl Transform. Version arXiv:2609.28993v1 (CC BY 4.0), distributed under the Creative Commons Attribution 4.0 International licence, http://creativecommons.org/licenses/by/4.0/, as recorded in the arXiv OAI licence metadata for this exact version and shown on the abstract page https://arxiv.org/abs/2609.28993v1. This item is an editorial re-presentation of the paper's uniform comparison of the Dunkl-weighted measures of two concentric cubes, which is the technical heart of the necessity half of its Logvinenko-Sereda theorem for the Dunkl transform, printed as Lemma 4.1. A full rights notice is delivered at licenses/arxiv-2609.28993-CC-BY-4.0.txt.
\noindent\textit{Editorial scope.} Counted unit: the paper's uniform comparison of the Dunkl-weighted measures of two concentric cubes, which is the technical heart of the necessity half of its Logvinenko-Sereda theorem for the Dunkl transform, printed as Lemma 4.1 (source0.tex lines 729--749 for the statement and 750--990 for the proof), delivered together with the paper's complete original proof of it as the single primary proof body. No mathematics is written, repaired or re-derived: statement and proof are the paper's own text line for line, in the paper's own notation. Every cross-reference inside the retained spans resolves inside the delivered document. Printed numbers are the paper's own: the wrapper restores the paper's counters before the retained material, so the counted statement prints as Lemma 4.1 and the two displays of the statement carry the paper's own equation numbers (24) and (25). Omitted from this package and not redistributed: the paper's preamble and front matter as such (of which only the title and the author names are reproduced in the wrapper); the sections, lemmas, theorems and examples that lie outside the retained spans; and the paper's bibliography with its bib data file. No citation command occurs in any retained span, so the delivered document carries no bibliography and no reference or citation command of the delivered text was rewritten or adapted. Printed slips kept literal, among them the paper's own wording and its own choice of symbols. Layout and class: the delivered wrapper is the lane's pinned amsart editorial wrapper at 10pt on A4, to which this package adds the paper's own theorem-like declarations and the shorthand macros its retained text uses, all copied verbatim from the paper's source. No publisher class file, no publisher style file and no upstream script is redistributed. Assets: no retained span of this package contains a figure environment or a graphics-inclusion command, no image or artwork file exists in or is redistributed with this package, and the manifest and the delivery evidence both carry an empty asset-dependency list. A full rights notice is delivered at licenses/arxiv-2609.28993-CC-BY-4.0.txt.
\setcounter{section}{4}
\setcounter{theorem}{0}
\setcounter{equation}{23}
\begin{lemma}\label{L4001}
Let
\[
Q=\biggl(-\frac12,\frac12\biggr)^d,\quad Q_0=x_0+LQ,\quad Q_1=x_0+(L+M)Q,
\]
where $M\ge 0$ is fixed. Then there exists a constant $C>0$ depending only on $d,R,k$ such that for all $x_0\in\mathbb{R}^d$ and $L>0$,
\begin{equation}\label{E4001}
0\le 1-\frac{\mu_k(Q_0)}{\mu_k(Q_1)}\le C\,\frac{M}{L+M}.
\end{equation}

Therefore,
\begin{equation}\label{E4002}
\lim_{L\to\infty}\sup_{x_0\in\mathbb{R}^d}\left|\frac{\mu_k(x_0+LQ)}{\mu_k(x_0+(L+M)Q)}-1\right|=0.
\end{equation}

In particular, for any $\delta>0$, there exists $K>0$ such that for all $L>K$, uniformly in $x_0\in\mathbb{R}^d$,
\[
\frac{\mu_k(Q_0)}{\mu_k(Q_1)}>1-\delta.
\]

\end{lemma}

\begin{proof}
% [R8-LANG-041] Improved precision, flow, or concision in the proof exposition.
The cases $M=0$ and $R_+=\emptyset$ are immediate. We therefore assume $M>0$ and $R_+\ne\emptyset$.
For Lebesgue measure, the volume ratio of two concentric cubes tends to $1$ as $L$ grows. To account for the weight, we show that the strip between the cubes carries only a small fraction of the total mass. The proof proceeds in three steps: we estimate the weight, control the strip, and then complete the argument.


\emph{Step 1. A one-dimensional weight estimate.}

 Fix $m\ge1$ and $\kappa\ge 0$. Let
\[
h(t)=A\prod_{j=1}^{r}|a_j t+b_j|^{\beta_j},
\]
where
\[
A\ge 0,\qquad r\le m,\qquad \beta_j\ge 0,\qquad \sum_{j=1}^{r}\beta_j\le \kappa.
\]
There is a constant depending only on $m$ and $\kappa$, which may be chosen as
\[
C_{m,\kappa}=2(1+4m)^{\kappa},
\]
such that for any finite interval $I\subset\mathbb{R}$ and any measurable set $E\subset I$, we have
\[
\int_{E} h(t)\,\mathrm{d}t \le C_{m,\kappa}\frac{|E|}{|I|}\int_{I} h(t)\,\mathrm{d}t.
\]

By an affine transformation, we may reduce the interval $I$ to $[0,1]$.
If a factor vanishes identically with a positive exponent, then $h\equiv0$ and there is nothing to prove. All nonzero constant factors corresponding to $a_j=0$ can be absorbed into $A$;
for $a_j\neq 0$, write
\[
|a_j t + b_j|^{\beta_j}=|a_j|^{\beta_j}|t-\rho_j|^{\beta_j},\qquad \rho_j=-\frac{b_j}{a_j}.
\]
Thus it suffices to consider
\[
h(t)=A_0\prod_{j=1}^{r}|t-\rho_j|^{\beta_j},\qquad t\in[0,1].
\]
Let
\[
\eta:=\frac{1}{4m},
\]
and define
\[
F:=[0,1]\setminus \bigcup_{j=1}^{r}(\rho_j-\eta,\rho_j+\eta).
\]
Since $r\le m$, we have
\begin{equation}
|F|\ge 1-2r\eta\ge 1-\frac12=\frac12.
\end{equation}

For any $t\in F$, we have
\[
|t-\rho_j|\ge \eta.
\]
Therefore, for any $s\in[0,1]$,
\[
\begin{aligned}
|s-\rho_j| &\le |s-t|+|t-\rho_j| \\
&\le 1+|t-\rho_j| \\
&\le (1+\eta^{-1})|t-\rho_j| \\
&=(1+4m)|t-\rho_j|.
\end{aligned}
\]
Hence
\[
h(s)\le (1+4m)^{\sum_j\beta_j}h(t)\le (1+4m)^{\kappa}h(t).
\]
Taking the supremum over $s\in[0,1]$, we obtain
\[
\|h\|_{L^{\infty}(0,1)}\le (1+4m)^{\kappa}h(t),\qquad t\in F.
\]
Integrating over $t\in F$, and using $|F|\ge 1/2$, we get
\[
\begin{aligned}
\int_{0}^{1}h(t)\,\mathrm{d}t &\ge \int_{F}h(t)\,\mathrm{d}t \\
&\ge \frac{|F|}{(1+4m)^{\kappa}}\|h\|_{L^{\infty}(0,1)} \\
&\ge \frac{1}{2(1+4m)^{\kappa}}\|h\|_{L^{\infty}(0,1)}.
\end{aligned}
\]
Therefore
\[
\|h\|_{L^{\infty}(0,1)}\le 2(1+4m)^{\kappa}\int_{0}^{1}h(t)\,\mathrm{d}t.
\]
Thus for any $E\subset [0,1]$,
\[
\int_{E} h(t)\,\mathrm{d}t \le |E|\,\|h\|_{L^{\infty}(0,1)}
\le 2(1+4m)^{\kappa}|E|\int_{0}^{1} h(t)\,\mathrm{d}t.
\]
Scaling back to the general interval $I$ via an affine transformation, we obtain
\[
\int_{E} h(t)\,\mathrm{d}t \le 2(1+4m)^{\kappa}\frac{|E|}{|I|}\int_{I} h(t)\,\mathrm{d}t.
\]
This proves the claim.

\emph{Step 2. A uniform estimate for the boundary strip.}

Fix
\[
a\in\mathbb{R}^d,\qquad \varepsilon>0,
\]
and define
\[
C_{\mathrm{in}}:=a+Q,\qquad C_{\mathrm{out}}:=a+(1+\varepsilon)Q.
\]
Since $Q=\bigl(-1/2,1/2\bigr)^d$, we have
\begin{equation}\label{EPL05}
C_{\mathrm{in}}\subset C_{\mathrm{out}}.
\end{equation}
For the $i$-th coordinate, let
\[
I_i:=\biggl(a_i-\frac{1+\varepsilon}{2},\,a_i+\frac{1+\varepsilon}{2}\biggr),
\]
and
\[
J_i:=\biggl(a_i-\frac12,\,a_i+\frac12\biggr).
\]
Then
\[
|I_i|=1+\varepsilon,\qquad |I_i\setminus J_i|=\varepsilon.
\]
Define the $i$-th boundary strip
\[
S_i:=\{x\in C_{\mathrm{out}}: x_i\in I_i\setminus J_i\}.
\]
Clearly
\begin{equation}
C_{\mathrm{out}}\setminus C_{\mathrm{in}}\subset \bigcup_{i=1}^{d} S_i.
\end{equation}
Fix all coordinates except the $i$-th one:
\[
\widehat{x}_i = (x_1,\dots,x_{i-1},x_{i+1},\dots,x_d).
\]
With the remaining coordinates fixed, write $w$ as a function of $t=x_i$:
\[
\begin{aligned}
h_{\widehat{x}_i}(t) &:= w(x_1,\dots,x_{i-1},t,x_{i+1},\dots,x_d) \\
&= \prod_{\alpha\in R_+} \biggl| \alpha_i t + \sum_{j\neq i}\alpha_j x_j \biggr|^{2k_\alpha}.
\end{aligned}
\]
Note that the number of factors is at most $|R_+|$ and the sum of all exponents is $2\sum_{\alpha\in R_+} k(\alpha)=\kappa$,

Applying the claim in Step 1 with $m=|R_+|$ and $C_* := 2\bigl(1+4|R_+|\bigr)^{\kappa}$, we obtain

\begin{equation}
\int_{I_i\setminus J_i} h_{\widehat{x}_i}(t)\,\mathrm{d}t
\le C_*\frac{\varepsilon}{1+\varepsilon}\int_{I_i} h_{\widehat{x}_i}(t)\,\mathrm{d}t.
\end{equation}
Integrating over the remaining $d-1$ coordinates and applying Fubini's theorem, we obtain
\begin{equation}
\mu_k(S_i)\le C_*\frac{\varepsilon}{1+\varepsilon}\,\mu_k(C_{\mathrm{out}}).
\end{equation}
Summing over $i=1,\dots,d$ and using \eqref{EPL05}, we obtain
\[
\begin{aligned}
\mu_k(C_{\mathrm{out}}\setminus C_{\mathrm{in}})
&\le \sum_{i=1}^{d}\mu_k(S_i) \\
&\le d C_*\frac{\varepsilon}{1+\varepsilon}\,\mu_k(C_{\mathrm{out}}).
\end{aligned}
\]
Therefore
\[
\begin{aligned}
1-\frac{\mu_k(C_{\mathrm{in}})}{\mu_k(C_{\mathrm{out}})}
&=\frac{\mu_k(C_{\mathrm{out}}\setminus C_{\mathrm{in}})}{\mu_k(C_{\mathrm{out}})} \\
&\le d C_*\frac{\varepsilon}{1+\varepsilon}.
\end{aligned}
\]

Hence for all $a\in\mathbb{R}^d$, uniformly,
\begin{equation}\label{EPL11}
\frac{\mu_k(a+Q)}{\mu_k\bigl(a+(1+\varepsilon)Q\bigr)}
\ge 1- d C_*\frac{\varepsilon}{1+\varepsilon}.
\end{equation}


\emph{Step 3. Rescaling by homogeneity.}

\[
w(\lambda x)=\lambda^{\kappa} w(x),\quad \lambda>0.
\]
Thus for any measurable set $A$,
\[
\mu_k(\lambda A)=\lambda^{d_k}\mu_k(A).
\]
Set
\[
a:=\frac{x_0}{L},\qquad \varepsilon:=\frac{M}{L}.
\]
Then
\[
Q_0=x_0+LQ=L(a+Q),
\]
and
\[
Q_1=x_0+(L+M)Q=L\bigl(a+(1+\varepsilon)Q\bigr).
\]
By homogeneity,
\[
\mu_k(Q_0)=L^{d_k}\mu_k(a+Q),
\]
and
\[
\mu_k(Q_1)=L^{d_k}\mu_k\bigl(a+(1+\varepsilon)Q\bigr).
\]
Therefore
\[
\frac{\mu_k(Q_0)}{\mu_k(Q_1)}
=\frac{\mu_k(a+Q)}{\mu_k\bigl(a+(1+\varepsilon)Q\bigr)}.
\]
Substituting \eqref{EPL11}, we obtain
\begin{equation}\label{EPL15}
\begin{aligned}
\frac{\mu_k(Q_0)}{\mu_k(Q_1)}
&\ge 1- d C_*\frac{M/L}{1+M/L} \\
&=1- d C_*\frac{M}{L+M}.
\end{aligned}
\end{equation}
We may therefore take
\[
C_{d,R,k}:= d C_* = 2d\bigl(1+4|R_+|\bigr)^{\kappa},
\]
so that
\[
0\le 1-\frac{\mu_k(Q_0)}{\mu_k(Q_1)}
\le C_{d,R,k}\frac{M}{L+M}.
\]
The constant in this estimate is independent of $x_0$.

Let $\delta>0$ be arbitrary. Set
\[
K:=\frac{C_{d,R,k}M}{\delta}.
\]
Whenever $L>K$,
\[
C_{d,R,k}\frac{M}{L+M}<C_{d,R,k}\frac{M}{L}<\delta.
\]
Thus by \eqref{EPL15}, for every $x_0\in\mathbb{R}^d$, we have
\[
\frac{\mu_k(x_0+LQ)}{\mu_k\bigl(x_0+(L+M)Q\bigr)}>1-\delta.
\]
Since $K$ is independent of $x_0$, the conclusion holds uniformly in $x_0$, giving \eqref{E4002}.

\end{proof}

\end{document}
